\newpage{}
\begin{abstract}
    The use of convolutional neural networks (CNN) to analyse patterns in transiting exoplanets has been observed as an upcoming method in filtering false positive signals. However, for best results, papers focus on preprocessing and cleaning the data heavily to make classification easier on the model. This paper focuses on improving the model to handle Kepler raw data for classification. Using lightly processed data, a binary and multiclass CNN were developed to classify light curves from Kepler as transit or shallow/deep eclipsing binaries. The binary model showed impressive results with \textbf{(How much)}\% accuracy score while the multiclass model had an accuracy of \textbf{(What score?)}\%. The main distinction comes from the tricky nature of shallow eclipsing binaries which are very common in observation. This research focuses on the developing process of using machine learning models to search for exoplanets with minimal manual processing requirements.
\end{abstract}

\section{Introduction}

Prior to 1992, the existence of planets beyond the Solar System was largely speculative. This changed with the first observational evidence of extrasolar objects in the early 1990s \citep{First_Exoplanet, First_Exoplanet_RV}. These confirmations transformed the field from theoretical speculation into an observational science. Since then, the rate of discovery has increased rapidly, and by 2008 over 263 exoplanet systems had been confirmed \citep{Exoplanet_numbers_2008}. As of August 2026, over 6000 exoplanets have been confirmed, and thousands more are candidates awaiting further observation \citep{christiansen2025nasaexoplanetarchiveexoplanet}.\footnote{Confirmed planet counts obtained from the NASA Exoplanet Archive Confirmed Planets Table, accessed March 2026: \url{https://exoplanetarchive.ipac.caltech.edu/}}

This rapid growth in detections has been driven by advances in observational techniques and instrumentation. 51 Pegasi b was the first exoplanet discovered orbiting a Sun-like star. It was detected using radial velocity (RV) measurements of its host star \citep{First_Exoplanet_RV}. As the number of candidates increased, manual and low-throughput detection approaches became inefficient. Although radial velocity measurements remain widely used, transit photometry has become the dominant detection method \citep{How_Many_Transits}.

When analysing transit data, a major source of false positive transit-like signals is eclipsing binary systems, where two stars orbit one another and periodically pass in front of each other from the observer's perspective \citep{2018ApJS..235...38T, Morton_2016}. These systems can produce repeated decreases in observed flux that resemble planetary transits, making them an important contaminant when identifying exoplanet candidates.

This paper trains machine learning models to classify light curves collected from NASA's \textit{Kepler} mission as either transiting exoplanets or eclipsing binaries. Since eclipsing binaries and other astrophysical systems can produce transit-like false positive signals, an automated classifier could reduce the need for manual vetting. Previous studies have explored machine learning approaches for false positive identification, often using heavily processed or phase-folded light curves. The goal of this work is to test whether a model can achieve useful classification performance using relatively limited preprocessing of the input data.

